% sections/01_introduction.tex
\chapter{Introduction}
\label{ch:introduction}

The generation of organic waste has become one of the defining environmental challenges of our time, particularly as urbanization and industrial activity continue to expand. A major contributor is food waste, a global stream that in the European Union alone is estimated at 88 million tonnes per year, excluding inedible parts removed from the supply chain \parencite{scherhaufer2018}. When considering only post-processing activities – wholesale, retail, food service and households – approximately 66.3 million tonnes of food waste are produced annually in the EU-27+1 \parencite{albizzati2021}. The scale of this loss highlights the urgent need for improved management strategies that move beyond simple disposal.

The environmental burden of food waste is substantial. It accounts for roughly 15–16\% of the total climate impact of the entire food supply chain, translating to about 186 Mt CO\textsubscript{2}-eq of greenhouse gas emissions, alongside significant contributions to acidification and eutrophication \parencite{scherhaufer2018}. The majority of these impacts arise during primary production, making prevention the single most effective strategy; however, when waste cannot be avoided, its valorisation through anaerobic digestion or composting is favoured over landfill disposal \parencite{albizzati2021, scherhaufer2018}. Despite this, much food waste continues to be landfilled or incinerated, forfeiting the opportunity to recover its embedded energy and nutrients.

Parallel to the food waste challenge, the rapid expansion of biodiesel production has generated an excess of another organic by-product: crude glycerol. Formed in the transesterification of triglycerides, approximately 10 kg of crude glycerol are produced for every 100 kg of biodiesel \parencite{garlapati2016}. Global biodiesel output, driven by renewable energy mandates, has caused the glycerol market to swell – with crude glycerol supply projected to reach about 6.33 million tonnes by 2025, of which nearly 680 000 tonnes are expected to be highly impure due to the shift towards waste-based feedstocks \parencite{attarbachi2023}.

Crude glycerol composition varies widely depending on the oil source and catalyst used, but it typically contains methanol, soaps, free fatty acids, salts and other matter organic non-glycerol (MONG) compounds \parencite{anand2012, attarbachi2023}. These impurities are known to inhibit microbial growth and fermentation, as demonstrated by the severe reduction in 1,3-propanediol production when untreated crude glycerol from jatropha, soybean or sunflower oils was used \parencite{anand2012}. Purification to technical or pharmaceutical grade is energy-intensive and often uneconomical for small and medium biodiesel plants \parencite{attarbachi2023}, prompting the search for direct, cost-effective valorisation routes. Because of its high organic content and ready biodegradability, crude glycerol has attracted attention as an ideal co-substrate for anaerobic digestion, where it can serve as a concentrated carbon source while being simultaneously detoxified by dilution and microbial synergy \parencite{gebreegziabher2025, tomczak2025}.

Anaerobic digestion (AD) is a mature biological process in which consortia of microorganisms convert organic matter into biogas – primarily methane (55–65\%) and carbon dioxide – under oxygen-free conditions, passing through the sequential stages of hydrolysis, acidogenesis, acetogenesis and methanogenesis \parencite{gebreegziabher2025}. AD offers a unique combination of benefits: it produces renewable, storable energy that can replace fossil fuels; generates a nutrient-rich digestate that can be used as organic fertiliser; and reduces the environmental load of organic wastes by avoiding landfilling and its associated methane emissions \parencite{scherhaufer2018, gebreegziabher2025}.

Within this framework, anaerobic co-digestion – the simultaneous treatment of two or more substrates – has emerged as a particularly effective strategy. Co-digesting crude glycerol with nitrogen-rich materials such as animal manure, sewage sludge or food waste balances the carbon-to-nitrogen ratio, dilutes potential inhibitors, provides essential macro- and micronutrients, and fosters synergistic microbial interactions \parencite{tomczak2025, gebreegziabher2025}. The result can be a dramatic improvement in biogas and methane yield: studies report increases ranging from 14\% to 226\% when crude glycerol is added to a variety of base substrates, with methane contents generally exceeding 60\% as long as glycerol does not exceed about 8\% (v/v) of the feed \parencite{gebreegziabher2025, tomczak2025}.

However, excessive glycerol loads can lead to volatile fatty acid accumulation, pH drop and methanogen inhibition, underscoring the need for careful optimisation \parencite{gebreegziabher2025, tomczak2025}. Overall, the convergence of abundant food waste and surplus crude glycerol supplies, together with the proven capability of anaerobic digestion to convert these streams into renewable methane and fertiliser, represents a tangible route towards a circular, low-carbon bioeconomy.

\chapter{Objectives}
\label{ch:objectives}

\section{General Objectives}
\label{sec:general_objectives}

The primary objective of this dissertation project is to evaluate the efficiency of a single-stage anaerobic co-digestion system using organic food waste and glycerol as substrates.The investigation will be conducted by determining biogas production under different operational conditions, with a focus on industrial feasibility.


The central hypothesis is that the controlled addition of glycerol, a by-product of the biodiesel industry, enhances biogas production due to its high energy content and rapid biodegradability, provided that adequate operational conditions are maintained to avoid volatile fatty acid overload.

\section{Specific Objectives}
\label{sec:specific_objectives}

\begin{enumerate}
    \item Characterize the substrates (food waste and glycerol), evaluating physicochemical parameters relevant to anaerobic digestion.
    \item Operate the single-stage anaerobic co-digestion system under different operational conditions, assessing process stability.
    \item Determine biogas production as a function of substrate proportions and operational conditions.
    \item Evaluate the techno-economic feasibility of using glycerol as a co-substrate, considering energy yield and potential for industrial application.
\end{enumerate}

